% GENERATED FILE. Edit canonical lessons, then run npm run generate:books.
% canonical-source-hash: fnv1a64:100d82b9cf72a4b0
% canonical-lessons: ES-C512-producirse, ES-C512-suceso, ES-C512-acontecimiento, ES-C512-acontecer, ES-C512-tener-lugar

\chapter{Lo que pasó}
\label{ch:es-lo-que-paso}
\hlchaptermodality{512}

\begin{chapteropening}
\textbf{By the end of this chapter:} \emph{I can say what happened: that something occurred, came about or took place, and call it an incident or an event.}

Complete the last lesson of chapter 512: I can say what happened: that something occurred, came about or took place, and call it an incident or an event.
\end{chapteropening}

\section[producirse]{producirse — to occur, to take place (of an event)}

\label{lesson:ES-C512-producirse}



\begin{quote}
Before the new one: say the Spanish for to seem, then the Spanish for to keep.
\end{quote}

\subsection*{You'll want to know: producirse}
\textbf{producirse} — \textquotedblleft{}to occur, to take place (of an event)\textquotedblright{}.

\begin{quote}
El accidente se produjo a las seis de la tarde.
\end{quote}

\begin{quote}
\emph{The accident occurred at six in the evening.}
\end{quote}

Common in the news. It changes like \emph{conducir}: \emph{se produjo}, \emph{se produjeron}.

\begin{cousinweb}
\emph{producirse} is built on \emph{producir} 'to produce', a learned borrowing from Latin \emph{prōdūcere} 'to lead forth, to bring out' (EtymDB), from \emph{prō} 'forward' and \emph{dūcere} 'to lead'. An event that occurs is one brought out into the open.
\end{cousinweb}

\begin{grammarlens}[title={What you've built}]
One more word for telling, explaining and describing at B1.
\end{grammarlens}

\subsection*{Your turn}
\begin{itemize}
\raggedright
  \item \emph{Say it:} \emph{producirse}
  \item \emph{Say it:} \emph{producirse}, once more
  \item \emph{Say it:} a sentence of your own with \emph{producirse}
  \item \emph{Recall:} say the Spanish for to seem, then the Spanish for to keep, then say \emph{producirse} again
\end{itemize}

\begin{tcolorbox}[breakable,colback=teal!4,colframe=teal!35!black,title={Before you move on}]
What does producirse mean? (\textquotedblleft{}To occur, to take place (of an event)\textquotedblright{}.) Say it once more.
\end{tcolorbox}
\section[el suceso]{el suceso — the happening, the incident}

\label{lesson:ES-C512-suceso}



\begin{quote}
Before the new one: say the Spanish for to keep, then the Spanish for to occur.
\end{quote}

\subsection*{You'll want to know: el suceso}
\textbf{el suceso} — \textquotedblleft{}the happening, the incident\textquotedblright{}.

\begin{quote}
El suceso ocurrió el lunes.
\end{quote}

\begin{quote}
\emph{The incident took place on Monday.}
\end{quote}

In a newspaper, \emph{la sección de sucesos} is the page of accidents and crimes. It is a false friend: success is \emph{el éxito}.

\begin{cousinweb}
\emph{suceso} comes from Latin \emph{successus} 'an advance, an outcome' (EtymDB), built on \emph{succēdere} 'to come up after', from \emph{sub} and \emph{cēdere} 'to go' (Lewis \& Short).
\end{cousinweb}

\begin{grammarlens}[title={What you've built}]
One more word for telling, explaining and describing at B1.
\end{grammarlens}

\subsection*{Your turn}
\begin{itemize}
\raggedright
  \item \emph{Say it:} \emph{el suceso}
  \item \emph{Say it:} \emph{el suceso}, once more
  \item \emph{Say it:} a sentence of your own with \emph{el suceso}
  \item \emph{Recall:} say the Spanish for to keep, then the Spanish for to occur, then say \emph{el suceso} again
\end{itemize}

\begin{tcolorbox}[breakable,colback=teal!4,colframe=teal!35!black,title={Before you move on}]
What does el suceso mean? (\textquotedblleft{}The happening, the incident\textquotedblright{}.) Say it once more.
\end{tcolorbox}
\section[el acontecimiento]{el acontecimiento — the event, the happening}

\label{lesson:ES-C512-acontecimiento}



\begin{quote}
Before the new one: say the Spanish for to occur, then the Spanish for the happening.
\end{quote}

\subsection*{You'll want to know: el acontecimiento}
\textbf{el acontecimiento} — \textquotedblleft{}the event, the happening\textquotedblright{}.

\begin{quote}
La boda de la princesa fue un gran acontecimiento.
\end{quote}

\begin{quote}
\emph{The princess's wedding was a big event.}
\end{quote}

Used for important or public events.

\begin{cousinweb}
\emph{acontecimiento} is built on \emph{acontecer} 'to come about' with the ending \emph{-miento} (EtymDB). \emph{Acontecer} goes back to Latin \emph{contingere} 'to touch, to befall' (Lewis \& Short).
\end{cousinweb}

\begin{grammarlens}[title={What you've built}]
One more word for telling, explaining and describing at B1.
\end{grammarlens}

\subsection*{Your turn}
\begin{itemize}
\raggedright
  \item \emph{Say it:} \emph{el acontecimiento}
  \item \emph{Say it:} \emph{el acontecimiento}, once more
  \item \emph{Say it:} a sentence of your own with \emph{el acontecimiento}
  \item \emph{Recall:} say the Spanish for to occur, then the Spanish for the happening, then say \emph{el acontecimiento} again
\end{itemize}

\begin{tcolorbox}[breakable,colback=teal!4,colframe=teal!35!black,title={Before you move on}]
What does el acontecimiento mean? (\textquotedblleft{}The event, the happening\textquotedblright{}.) Say it once more.
\end{tcolorbox}
\section[acontecer]{acontecer — to come about; to happen}

\label{lesson:ES-C512-acontecer}



\begin{quote}
Before the new one: say the Spanish for the happening, then the Spanish for the event.
\end{quote}

\subsection*{You'll want to know: acontecer}
\textbf{acontecer} — \textquotedblleft{}to come about; to happen\textquotedblright{}.

\begin{quote}
Nadie sabe qué va a acontecer.
\end{quote}

\begin{quote}
\emph{Nobody knows what is going to come about.}
\end{quote}

A formal word, mostly in the third person: \emph{aconteció}. It takes a \emph{z} in \emph{acontezca}.

\begin{cousinweb}
\emph{acontecer} is built on old Spanish \emph{contecer} with the prefix \emph{a-} (EtymDB). \emph{Contecer} comes from spoken Latin \emph{*contingescere} (EtymDB), from \emph{contingere} 'to touch, to befall' (Lewis \& Short).
\end{cousinweb}

\begin{grammarlens}[title={What you've built}]
One more word for telling, explaining and describing at B1.
\end{grammarlens}

\subsection*{Your turn}
\begin{itemize}
\raggedright
  \item \emph{Say it:} \emph{acontecer}
  \item \emph{Say it:} \emph{acontecer}, once more
  \item \emph{Say it:} a sentence of your own with \emph{acontecer}
  \item \emph{Recall:} say the Spanish for the happening, then the Spanish for the event, then say \emph{acontecer} again
\end{itemize}

\begin{tcolorbox}[breakable,colback=teal!4,colframe=teal!35!black,title={Before you move on}]
What does acontecer mean? (\textquotedblleft{}To come about; to happen\textquotedblright{}.) Say it once more.
\end{tcolorbox}
\section[tener lugar]{tener lugar — to take place}

\label{lesson:ES-C512-tener-lugar}



\begin{quote}
Before the new one: say the Spanish for the event, then the Spanish for to come about.
\end{quote}

\subsection*{You'll want to know: tener lugar}
\textbf{tener lugar} — \textquotedblleft{}to take place\textquotedblright{}.

\begin{quote}
La reunión tuvo lugar en el ayuntamiento.
\end{quote}

\begin{quote}
\emph{The meeting took place at the town hall.}
\end{quote}

Used for planned events: \emph{La boda tendrá lugar en mayo.}

\begin{cousinweb}
\emph{tener lugar} is built on \emph{tener}, inherited from Latin \emph{tenēre} 'to hold' (EtymDB), and \emph{lugar}, which is usually traced to Latin \emph{locālis} 'of a place', from \emph{locus} 'a place' (Lewis \& Short). To take place is, word for word, to hold a place.
\end{cousinweb}

\begin{grammarlens}[title={What you've built}]
One more word for telling, explaining and describing at B1.
\end{grammarlens}

\subsection*{Your turn}
\begin{itemize}
\raggedright
  \item \emph{Say it:} \emph{tener lugar}
  \item \emph{Say it:} \emph{tener lugar}, once more
  \item \emph{Say it:} a sentence of your own with \emph{tener lugar}
  \item \emph{Recall:} say the Spanish for the event, then the Spanish for to come about, then say \emph{tener lugar} again
\end{itemize}

\begin{tcolorbox}[breakable,colback=teal!4,colframe=teal!35!black,title={Before you move on}]
What does tener lugar mean? (\textquotedblleft{}To take place\textquotedblright{}.) Say it once more.
\end{tcolorbox}
